\subsection{Pilot study (study01)}
The pilot used the curated handoff corpus: 1{,}232 AMP-labeled peptide sequences with synthetic FASTA
identifiers and unverified source provenance (8--60 aa) and 1{,}500 length-matched non-AMP negatives, evaluated
by 5-fold cross-validation with bootstrap intervals ($B=500$) and
label-permutation tests ($B=500$). All four models beat chance decisively
(Table~\ref{tab:pilot}). PepCNN led at AUROC $0.947$ [0.936, 0.955], with the
random forest on dipeptide composition surprisingly strong at $0.935$; the
GNN and logistic baseline trailed near $0.80$. The pilot shows an internal classifier difference on this corpus; because
the handoff FASTA identifiers are synthetic, we cannot independently
verify the provenance or annotation status of every source sequence.

\subsection{Scale-up with homology-aware splits (study07)}
The scale-up used the full fetched corpus: \Npos{} positives and \Nneg{}
length-matched negatives, clustered into \Nclusters{} shared-10-mer
components and split 80/10/10 by cluster (Section 3.3). Models trained on a
30{,}000-sequence subsample of the 40{,}776-sequence train split (8 epochs,
batch 128, weighted BCE). On the held-out test split of \Ntest{} sequences
(Table~\ref{tab:scaleup}), PepCNN attained AUROC $0.921$ [0.913, 0.929],
AUPRC $0.879$ --- clearly ahead of the random forest ($0.879$), the GNN
($0.877$), and logistic regression ($0.871$); all permutation $p=0.0020$.
The ordering inverts part of the pilot: with leakage controlled and $20\times$
more data, the learned-representation model separates from the feature
baselines, while the random forest's pilot advantage evaporates. This ordering is descriptive within the confounded scale-up. Because source/review status tracks class and multiple design choices change from the pilot, it does not establish that the CNN learns antimicrobial motifs better as data grow.

\subsection{APD6 positive-set comparison (study09)}
APD6 supplies positives curated separately from the UniProt keyword query, but the comparison still reuses reviewed UniProt negatives and a model trained on source-confounded labels. It cannot isolate cross-database biological transfer. Of 3{,}306 entries, 3{,}207 passed standard-residue and
length filters, and 1{,}412 survived decontamination against the training
split by the shared-10-mer criterion. Scored against the 2{,}925 held-out
test non-AMPs (Table~\ref{tab:apd}), PepCNN attained AUROC
$0.926$ [0.918, 0.934] on that selected comparison --- while the ensemble reached $0.909$ and the GNN $0.797$.
Sensitivity on decontaminated APD6 at a 5\% false-positive rate was $0.601$.
The limited reading: the CNN ranks many APD6 positives above the reviewed
UniProt test negatives, but the training-source confound still applies. Roughly two in five APD6 positives fall below the stringent operating threshold. The composition of those missed positives is unknown until an accession-grounded per-family analysis is run.

\subsection{Annealed design (studies 1 and 7)}
Running the simulated-annealing designer against the
pilot ensemble yielded ten candidates scoring $0.986$--$0.993$ with 3-mer
novelty $0.90$--$0.94$; against the scale-up ensemble, 24 anneals scored
$0.877$--$0.910$ (top 10 in Table~\ref{tab:designs}). The scale-up designs
score lower in absolute terms, but the difference in training data and model
calibration means this contrast does not establish improved calibration. All designs are \emph{in silico} hypotheses: no
wet-lab or even secondary-structure validation is claimed.

\subsection{Discovery screen (study08)}
We screened 852 reviewed UniProt entries of length 10--60 annotated only as
``uncharacterized'' or ``hypothetical'' protein (KW-0929 excluded) with the
study07 ensemble. The top 25 (top 15 in Table~\ref{tab:screen}) are enriched
for short, cysteine-rich, signal-peptide-like sequences. The top hit,
P0ACW4/P0ACW5 (identical 57-aa sequences from two accessions, ensemble
$0.908$, pairwise 3-mer novelty $0.95$ against the training reference), is a small Cys-rich protein; such sequences
superficially resemble defensin-like AMPs, which is plausibly why the model
scores them highly --- and equally plausibly reflects a learned prior for
Cys-rich small proteins rather than specific antimicrobial function. We
report them as ranked candidates for downstream validation, nothing more.
